\section{Example: Magnetic Field Errors} \label{ex8}

The quality of an undulator depends essentially on the errors of the
magnetic field. Field errors not only effect the trajectory and hence e.g. the
orbit of electrons in  storage rings, but also reduce the spectral
brightness of the device, in particular for the higher harmonics.

\W~offers some options to study these effects. We enter:

\cboxs{\senv \grau \\
python3 ../python/wave\_example\_8.py\\
cp \$WAVE/input/wave\_example\_8.in wave.in\\
wave\\
wplot \\
}

We look at the spectrum of the flux-density. Click:

\cboxs{Plot / Spectra / Flux-density / Fd}

\begin{figure}[htb]
    \centering
    \includegraphics[width=0.8\textwidth]
    {/home/scheer/wav/doc/pp/wave_example_8_no-errors.pdf}
    \caption{Nineth harmonic of a planar undulator}
    \label{ex8fig1}
\end{figure}

The main errors of permanent magnet undulators arise from fluctuations
of the magnetization of the block. This can be a variation in the
strength or the direction of the blocks. In \wi~we set\\

\cboxss{BCRAN=0.001\\ UBANGERR(1)=0.2}

In addition, feel free to modify the start of the random generator  by changing \scap{IRNMODE} in the
namelist \scap{\$RANDOMN}. Of course, then you will reproduce the
plots of this example{\ldots}

Then we run \W again and look at the trajectory (Fig. \ref{ex8fig2})
and the spectrum by plotting the overview.\\

\begin{figure}[htb]
    \centering
    \includegraphics[width=0.8\textwidth]
    {/home/scheer/wav/doc/pp/wave_example_8_errors.pdf}
    \caption{Horizontal, vertical trajectory , the nineth
    harmonic of an undulator with field errors}
    \label{ex8fig2}
\end{figure}

The field errors result in a spoilt trajectory with a remaining kick
(to see it better, keep the left mouse button pressed on the 3d
trajectory plot and move the mouse to rotate the plot).

In addition the on-axis spectrum of the flux-density is ruined (Fig.
\ref{ex8fig2}). But the later is mainly a consequence of the stearing of
the undulator, i.e. we can move the position of the observer to take
this into account by looking for the point the trajectory points to:\\

\cboxss{OBS1Y=-8888.\\OBS1Z=-8888.\\}

\begin{figure}[htb]
    \centering
    \includegraphics[width=0.8\textwidth]
    {/home/scheer/wav/doc/pp/wave_example_8_fd_obs.pdf}
    \caption{Nineth harmonic at the shifted observation point}
    \label{ex8fig3}
\end{figure}

This triggers a straight line fit of the trajectory in order to
estimate the best obervation point. We run \W~and find the spectrum of
Fig. \ref{ex8fig3}. The peak of the nineth harmonic is now higher; however, a
significant reduction of the flux-density and a distorted shape remains due to
the field errors.\\

We have used Gaussian errors in this example so far. Normally, one
would not accept magnets with errors above a certain limit, let's say
one standard deviation. This can be done by setting

\cboxss{BCRANSIG=1.0\\UBANSIG(1)=1.0\\}

Run \W and enter in the plotting shell:

\cboxss{~\\[-10pt]
zone(2,2)\\
nz()\\
zone(2,2,3,'same')\\
ny()\\
zone(2,1,2,'same')\\
hcfluxden()
}
%\newpage

The result is shown in Fig. \ref{ex8fig5}. We see a ''realistic''
trajectory and a good nineth harmonic.

These options allow the user to estimate the required quality of the
magnets and trajectory of the undulator. There is another option
\scap{USIGOFFY(1)}, which limits the random walk of the trajectory due
to angle errors of the magnets. This simulates the shimming procedure
of the undulator and also allows to estimate the required quality of
the trajectory.

However, these tools are provided to produce
''realistic'' trajectories and give hints about the impact of field
errors. They do probably not replace a thoroughly analysis of the
field measurement, analysis, and shimming of an undulator.

\begin{figure}[htb]
    \centering
    %\framebox(13,16){
          \includegraphics[width=0.7\textwidth]
	  {/home/scheer/wav/doc/pp/wave_example_8_lim_errors.pdf}
    %} % end framebox
    \caption{Hori. and vert. trajectories  and spectral flux-density
    of an undulator with limited errors}
    \label{ex8fig5}
\end{figure}

%\newpage
\underline{Proposed exercises:}

Check the option \scap{NURANMOD}. With \scap{NURANMOD=2}, the error
\scap{BCRAN} can be applied to neighboring poles. This simulates a
perfect shimming of the trajectory, while keeping the phase errors and
there influence on the spectrum.

Dump the magnets of the model using \scap{IKRESTOR}, change some magnet
properties, and try to use the file to investigate the effects.\\

Check the option \scap{IBERROR} and the namelist \scap{\$BERRORN} and
try to use it.\\
